\section{实验}

\subsection{设置}

\subsubsection{任务}
我们既在标准 LLM 基准上开展实验，也在更复杂的智能体任务上进行实验。任务、数据集和评估指标如下。

\begin{table}[htb]
  \centering
  \caption{\textbf{标准 LLM 基准上的结果。}HotPotQA 单元格依次报告精确匹配/F1；MATH 报告准确率；HumanEval 报告 Pass@1。}
  \label{tab:benchmark_results}
  \small
  \renewcommand\tabcolsep{6pt}
  \renewcommand\arraystretch{1.16}
  \begin{tabular}{l|cc|cc|cc}
  \Xhline{1.5pt}
  \textbf{方法} & \multicolumn{2}{c|}{\textbf{HotPotQA}} & \multicolumn{2}{c|}{\textbf{MATH}} & \multicolumn{2}{c}{\textbf{HumanEval}} \\
                   & \textbf{GPT-3.5} & \textbf{GPT-4} & \textbf{GPT-3.5} & \textbf{GPT-4} & \textbf{GPT-3.5} & \textbf{GPT-4} \\
  \Xhline{1.5pt}
  \bfseries GPTs & 24 / 38.8 & 33 / 44.3 & 23.2 & 53.1 & 59.2 & 71.7 \\
  \bfseries Agents & 27 / 37.5 & 39 / 49.8 & 23.8 & 56.0 & 59.5 & 85.0 \\
  \bfseries Agents w/ AutoPE & 29 / 39.8 & 38 / 50.3 & 22.5 & 57.2 & 63.5 & 82.3 \\
  \bfseries DSPy & 35 / 43.9 & 40 / 50.5 & 17.3 & 48.4 & \bfseries 66.7 & 77.3 \\ \hline
  \bfseries 本文方法 & \bfseries 35 / 44.8 & \bfseries 41 / 54.0 & \bfseries 38.8 & \bfseries 60.7 & 64.5 & \bfseries 85.8 \\
  \Xhline{1.5pt}
  \end{tabular}
\end{table}


\paragraph{标准基准}
我们在 HotpotQA~\citep{yang-etal-2018-hotpotqa}、MATH~\citep{hendrycksmath2021}和 HumanEval~\citep{chen2021codex}等标准 LLM 基准上开展实验。HotPotQA 是一项多跳问答任务，对丰富的背景知识提出挑战。由于我们发现数据集中的“hard”划分对语言智能体更具挑战性，因此使用该划分。MATH 收录了具有挑战性的数学竞赛题。HumanEval 是一个评估集，要求 LLM 或智能体根据文档字符串合成程序。评估指标方面，HotPotQA 使用 F1 和精确匹配，MATH 使用准确率，HumanEval 使用 Pass@1。我们在这些数据集上禁用工具，以保证结果能够与既有文献进行有意义的比较。

\paragraph{复杂智能体任务}
我们把\textbf{创意写作}和\textbf{软件开发}视为两项复杂智能体任务。对于创意写作任务，我们遵循\citet{yao2023tree}的设置，向智能体提供 4 个随机句子，要求它们写出一篇连贯的 4 段文章，并让 4 个段落分别以这 4 个输入句子结尾。该任务具有开放性和探索性，既考验创造性思维，也考验高层规划。我们遵循\citet{yao2023tree}，使用 GPT-4 评分来评估文章。另一方面，软件开发任务要求智能体系统根据一份简单的产品需求文档（PRD）开发出\textit{可执行}软件。我们依据所生成软件的\textit{可执行性}评估各对比智能体；该指标用 $1$ 到 $4$ 的数值分数衡量，对应逐级提高的执行能力。具体来说，$1$ 分表示执行失败，$2$ 分表示代码能够成功运行，$3$ 分表示符合预期工作流程，$4$ 分表示与预期完全一致。

\subsubsection{基线}

我们将所提出的方法与以下基线进行比较：
\begin{itemize}
    \item \textbf{GPTs}：一种简单基线，使用 GPT 和精心设计的提示；
    \item \textbf{Agents}：使用 Agents~\citep{zhou2023agents}框架\footnote{我们也测试了 OpenAgents 和 AgentVerse 等其他智能体框架，并得到相似结果。}实现的语言智能体方法，配有精心设计的提示、工具和流水线；
    \item \textbf{DSPy}：一种 LLM 流水线优化框架，能够搜索提示组件的最佳组合。对于评估指标无法用方程和代码定义的复杂智能体任务，它并不适用；
    \item \textbf{Agents + AutoPE}：一个变体，其中智能体流水线各节点的提示由 LLM 按照\citet{yang2024large}所述方法进行优化。该变体不涉及语言梯度反向传播，也不涉及基于语言梯度的优化。
\end{itemize}

我们同时报告 GPT-3.5 和 GPT-4 的结果。GPT-3.5 使用 \texttt{gpt-3.5-turbo-0125} 端点，GPT-4 使用 \texttt{gpt-4-turbo-0409} 端点。对于本文方法，我们以\textbf{Agents}基线为起点，再在其上开展\baby{}。

\subsection{结果}

\begin{table}[htb]
  \centering
  \begin{minipage}[t]{.53\linewidth}
      \centering
      \caption{软件开发任务结果。}
      \label{tab:software}
      \small
      \renewcommand\tabcolsep{4pt}
      \begin{tabular}{l|ccc}
      \Xhline{1.5pt}
      \textbf{任务} & \textbf{GPTs} & \textbf{Agents} & \textbf{本文方法} \\
      \hline
      \bfseries Flappy Bird & 2 & 2 & 3 \\
      \bfseries 坦克大战 & 1 & 2 & 4 \\
      \bfseries 2048 & 1 & 2 & 4 \\
      \bfseries 贪吃蛇 & 2 & 3 & 4 \\
      \bfseries 打砖块 & 2 & 3 & 4 \\ \hline
      \bfseries 平均分 & 1.6 & 2.4 & \bfseries 3.8 \\
      \Xhline{1.5pt}
      \end{tabular}
  \end{minipage}%
  \hfill
  \begin{minipage}[t]{.43\linewidth}
      \centering
      \caption{创意写作任务结果。}
      \label{tab:writing}
      \small
      \renewcommand\tabcolsep{4pt}
      \begin{tabular}{l|cc}
      \Xhline{1.5pt}
      \textbf{方法} & \textbf{GPT-3.5} & \textbf{GPT-4} \\
      \hline
      \bfseries GPTs & 4.0 & 6.0 \\
      \bfseries Agents & 4.2 & 6.0 \\
      \bfseries Agents w/ AutoPE & 4.4 & 6.5 \\
      \bfseries ToT & 3.8 & 6.8 \\ \hline
      \bfseries 本文方法 & \bfseries 6.9 & \bfseries 7.4 \\
      \Xhline{1.5pt}
      \end{tabular}
  \end{minipage}
\end{table}


\paragraph{LLM 基准上的结果}
标准 LLM 基准上的结果见表~\ref{tab:benchmark_results}。可以看到，所提出的\baby{}框架相对于所有对比方法均取得了一致提升；其中，在竞赛级基准 MATH 上的性能增益尤其显著。相比之下，传统的基于 LLM 的提示优化方法（Agents w/ AutoPE）与基于搜索的提示优化方法（DSPy）并不稳定：它们在某些情形下带来良好提升，在另一些情形下却导致性能显著下降。这表明，\baby{}框架更加鲁棒，能够更有效地优化语言智能体的整体性能。

\paragraph{复杂任务上的结果}
软件开发和创意写作的结果分别见表~\ref{tab:software}和表~\ref{tab:writing}。可以看到，本文方法在两项任务上均显著优于所有对比基线，而且性能差距甚至大于传统 LLM 基准上的差距。有趣的是，在创意写作任务上，本文方法甚至优于思维树（tree-of-thought）；后者是一种经过精心设计的提示工程与推理算法。我们发现，本文方法成功找到了规划、写作和修订流水线，并且各步骤的提示都得到了良好优化。我们还发现，\baby{}框架恢复出与 MetaGPT~\citep{hong2023metagpt}中专门面向软件开发而设计的标准操作流程相似的流程。这验证了所提出\baby{}框架在现实任务上的有效性：这些任务没有真实标签，整体性能也无法用方程或代码计算，因而不同于 DSPy 等基于搜索的算法所适用的情形。

\subsection{案例研究与分析}

\begin{figure}[tp]
    \centering
    \includegraphics[width=0.62\linewidth]{figs/case.pdf}
    \caption{在创意写作任务上开展的一项案例研究。}
    \label{fig:case}
\end{figure}

接下来，图~\ref{fig:case}给出了\baby{}框架优化动态的案例研究。可以看到，本文方法能够像人类专家开发语言智能体那样，有效开展提示工程并设计智能体流水线。

此外，我们发现，与神经网络初始化对训练十分重要类似，智能体系统的初始化也会对最终性能产生不可忽视的影响。总体上，我们发现，以最简单的方式初始化智能体、再让符号优化器执行优化通常是有帮助的。相反，如果初始智能体系统经过度工程化，性能往往会变得不稳定。由这一观察自然引出的扩展方向是：或许可以先在大规模、多样化任务上开展某种预训练，为通用智能体得到一种用途广泛的初始化，再用\baby{}将其适配到专门任务。我们还发现，与由准确率或 F1 评估性能的标准基准相比，本文方法在复杂现实任务上的成功更加显著且稳定。这表明，未来的智能体学习研究应当更多聚焦现实任务；智能体研究社区也应致力于构建侧重智能体学习评估、由多样化复杂智能体任务组成的基准，并研究稳健衡量进展的方法。
